\documentclass[10pt,a4paper]{article}
\usepackage[T1]{fontenc}
\usepackage[utf8]{inputenc}
\usepackage[ngerman]{babel}
\usepackage[a4paper,left=12mm,right=12mm,top=12mm,bottom=9mm]{geometry}
\usepackage{lmodern}
\usepackage[table]{xcolor}
\usepackage{tabularx}
\usepackage{array}
\usepackage{microtype}
\usepackage[most]{tcolorbox}
\usepackage{tikz}
\usetikzlibrary{arrows.meta,calc,positioning}
\usepackage{pifont}
\usepackage[normalem]{ulem}
\pagestyle{empty}
\setlength{\parindent}{0pt}
\setlength{\parskip}{0pt}
\setlength{\tabcolsep}{3pt}
\renewcommand{\familydefault}{\sfdefault}
\linespread{1.0}\selectfont
\renewcommand{\arraystretch}{1.02}
\definecolor{Navy}{HTML}{17324D}
\definecolor{Blue}{HTML}{3B6EA5}
\definecolor{Sky}{HTML}{EAF3FA}
\definecolor{Mint}{HTML}{E8F5EF}
\definecolor{Gold}{HTML}{F3B33D}
\definecolor{Ink}{HTML}{1E2935}
\definecolor{Muted}{HTML}{617080}
\definecolor{Line}{HTML}{D8E1E8}
\definecolor{CodeBg}{HTML}{F5F7F9}
% Gruppenfarben: dezente Abtönungen der vorhandenen Palette
\colorlet{GBer}{Blue!22}
\colorlet{GMue}{Gold!38}
\definecolor{GSch}{HTML}{CFE8DA}
\colorlet{GUa}{Blue!22}
\colorlet{GUb}{Gold!38}
\color{Ink}
\newcommand{\sectiontitle}[1]{%
  \vspace{.3mm}{\color{Blue}\bfseries\large #1}\par
  \vspace{.2mm}{\color{Blue}\rule{\linewidth}{.65pt}}\par\vspace{.3mm}}
\newcommand{\sqlkw}[1]{\textcolor{Blue}{\bfseries #1}}
\newcommand{\sqlbox}[1]{%
  \begin{tcolorbox}[colback=CodeBg,colframe=Line,arc=1.3mm,boxrule=.55pt,
    left=2.2mm,right=2.2mm,top=.5mm,bottom=.5mm,boxsep=0pt]
  {\ttfamily\fontsize{9.4}{11.2}\selectfont #1}
  \end{tcolorbox}}
\newcommand{\workbox}[1]{%
  \begin{tcolorbox}[colback=Sky,colframe=Blue,arc=1.2mm,boxrule=.55pt,
    left=2mm,right=2mm,top=.35mm,bottom=.35mm,boxsep=0pt]
  {\fontsize{9.2}{11}\selectfont\hyphenpenalty=10000\exhyphenpenalty=10000\textcolor{Blue}{\bfseries Arbeitsauftrag:} #1}
  \end{tcolorbox}}
\newcommand{\pageheader}{%
  \colorbox{Navy}{\parbox{\dimexpr\linewidth-2\fboxsep\relax}{%
    \vspace{1.5mm}\color{white}\begin{tabularx}{\linewidth}{@{}Xr@{}}
    {\small\bfseries INFORMATIK \textperiodcentered{} KLASSE 9} &
    \colorbox{Gold}{\color{Navy}\small\bfseries\strut LÖSUNGEN}
    \end{tabularx}
    \vspace{1mm}{\fontsize{17}{20}\selectfont\bfseries
    Aufgabe 6 -- Gruppierungen mit Group BY}\par
    \vspace{.6mm}{\large Sicherungsblatt}\vspace{1.5mm}}}}
\newcommand{\pagefooter}[1]{%
  \vfill{\color{Line}\rule{\linewidth}{.5pt}}\par\vspace{1mm}
  \begin{tabularx}{\linewidth}{@{}Xc r@{}}
  {\small\color{Muted}Klasse 9 \textperiodcentered{} Datenbanken} &
  {\small\color{Muted}#1 / 1} &
  {\small\color{Muted}Sicherungsblatt \textperiodcentered{} Aufgabe 6}
  \end{tabularx}}
\newcommand{\merkebox}[1]{%
  \begin{tcolorbox}[colback=Mint,colframe=Blue,arc=2mm,boxrule=.8pt,
    left=3mm,right=3mm,top=1.5mm,bottom=1.5mm]
  \textcolor{Blue}{\bfseries Merke}\par\vspace{.5mm}
  {\fontsize{9.6}{11.6}\selectfont #1}
  \end{tcolorbox}}
\newcommand{\smallcaption}[1]{{\small\color{Muted}#1}\par\vspace{.5mm}}
\newcommand{\miniresult}[2]{%
  {\fontsize{8.6}{10.2}\selectfont
  \begin{tabular}{|>{\ttfamily\bfseries}l|}\hline
  \rowcolor{Sky}#1\\\hline
  #2\\\hline
  \end{tabular}}}
% Bausteine der Gruppengrafik
\newcommand{\tabfont}{\fontsize{8.6}{10.2}\selectfont}
\tikzset{
  bigarrow/.style={-{Triangle[length=2.6mm,width=3.6mm]},line width=1.6pt,draw=Blue},
  member/.style={draw=Ink!45,line width=.4pt,rounded corners=.8mm,minimum height=4.2mm,
    inner sep=.6mm,font=\fontsize{8}{9.4}\selectfont,anchor=south},
  countbadge/.style={font=\fontsize{8.6}{10}\selectfont\ttfamily\bfseries,text=Blue,anchor=north}
}
% \pot{x links}{y oben}{Breite}{Höhe}{Farbe}{Aufschrift} -- Maße in mm (x=y=1mm)
\newcommand{\pot}[6]{%
  \fill[#5!45] (#1,#2) -- (#1,{#2-#4+2.5}) arc[start angle=180,end angle=270,radius=2.5mm]
    -- ({#1+#3-2.5},{#2-#4}) arc[start angle=270,end angle=360,radius=2.5mm] -- ({#1+#3},#2) -- cycle;
  \draw[Navy,line width=1.1pt,line join=round] ({#1-1.4},{#2+.6}) -- (#1,#2) -- (#1,{#2-#4+2.5})
    arc[start angle=180,end angle=270,radius=2.5mm] -- ({#1+#3-2.5},{#2-#4})
    arc[start angle=270,end angle=360,radius=2.5mm] -- ({#1+#3},#2) -- ({#1+#3+1.4},{#2+.6});
  \node[font=\fontsize{8.4}{10}\selectfont\ttfamily,anchor=south] at ({#1+#3/2},{#2+.8}) {#6};}
\newcommand{\graynote}[1]{{\fontsize{7.8}{9.2}\selectfont\color{Muted}#1}}
\begin{document}
\pageheader\par
\tcbset{before skip=.6mm,after skip=.6mm}
\sectiontitle{1 \textperiodcentered{} GROUP BY bildet für jeden Wert eine Gruppe}
\begin{minipage}[t]{.34\linewidth}\vspace{0pt}
\workbox{Wie viele Mitglieder stammen aus jeder Stadt?}
\end{minipage}\hfill
\begin{minipage}[t]{.33\linewidth}\vspace{0pt}
\sqlbox{\sqlkw{SELECT} city, \sqlkw{COUNT}(*) \sqlkw{AS} Anzahl\par\sqlkw{FROM} users\par\sqlkw{GROUP BY} city}
\end{minipage}\hfill
\begin{minipage}[t]{.29\linewidth}\vspace{0pt}
\smallcaption{Vergleich im Ausschnitt}
{\tabfont\begin{tabular}{|l|>{\ttfamily}l|}\hline
\rowcolor{Sky}ohne GROUP BY & COUNT(*) = 6\\\hline
mit GROUP BY city & 3 \textperiodcentered{} 2 \textperiodcentered{} 1\\\hline
\end{tabular}}
\end{minipage}\par\vspace{1.2mm}
\begin{tikzpicture}[x=1mm,y=1mm]
  % Ausschnitt aus users
  \node[anchor=north west,inner sep=0] at (0,0) {%
    \begin{minipage}{47mm}
    \smallcaption{Ausschnitt aus users (6 Zeilen)}
    {\tabfont\begin{tabular}{|>{\ttfamily\bfseries}c|l|l|}\hline
    \rowcolor{Sky}id & username & city\\\hline
    \rowcolor{GBer}4 & koch\_kicker & Berlin\\\hline
    \rowcolor{GMue}2 & wanderwut\_finn & München\\\hline
    \rowcolor{GSch}1 & zornigerfotograf & Schwerin\\\hline
    \rowcolor{GBer}34 & gartenlinse & Berlin\\\hline
    \rowcolor{GMue}11 & kunstnews\_luis & München\\\hline
    \rowcolor{GBer}37 & musikriese & Berlin\\\hline
    \end{tabular}}
    \end{minipage}};
  % Pfeil 1
  \draw[bigarrow] (48.8,-17) -- (57.2,-17);
  \node[font=\fontsize{8.4}{10}\selectfont\ttfamily\bfseries,text=Blue,align=center,anchor=south] at (52.6,-15.6) {GROUP BY\\city};
  \node[font=\fontsize{7.4}{8.6}\selectfont,text=Muted,align=center,anchor=north] at (52.8,-19) {gleicher\\city-Wert\\$\rightarrow$ gleicher\\Topf};
  % Töpfe (Gruppen)
  \pot{60}{-7}{28}{18.5}{GBer}{city = 'Berlin'}
  \node[member,fill=GBer] at (74,-24.3) {37 \textperiodcentered{} musikriese};
  \node[member,fill=GBer] at (74,-19.6) {34 \textperiodcentered{} gartenlinse};
  \node[member,fill=GBer] at (74,-14.9) {4 \textperiodcentered{} koch\_kicker};
  \node[countbadge] at (74,-27.2) {COUNT(*) = 3};
  \pot{90.5}{-7}{28}{18.5}{GMue}{city = 'München'}
  \node[member,fill=GMue] at (104.5,-24.3) {11 \textperiodcentered{} kunstnews\_luis};
  \node[member,fill=GMue] at (104.5,-19.6) {2 \textperiodcentered{} wanderwut\_finn};
  \node[countbadge] at (104.5,-27.2) {COUNT(*) = 2};
  \pot{121}{-7}{28}{18.5}{GSch}{city = 'Schwerin'}
  \node[member,fill=GSch] at (135,-24.3) {1 \textperiodcentered{} zornigerfotograf};
  \node[countbadge] at (135,-27.2) {COUNT(*) = 1};
  % Pfeil 2
  \draw[bigarrow] (151.5,-17) -- (157.5,-17);
  \node[font=\fontsize{7.4}{8.6}\selectfont,text=Muted,align=center,anchor=north] at (154.5,-19) {1 Zeile\\je\\Gruppe};
  % Ergebnisrelation
  \node[anchor=north west,inner sep=0] at (159,-6.2) {%
    \begin{minipage}{26.5mm}
    \smallcaption{Ergebnis}
    {\tabfont\begin{tabular}{|l|>{\ttfamily\bfseries}r|}\hline
    \rowcolor{Sky}city & Anzahl\\\hline
    \rowcolor{GBer}Berlin & 3\\\hline
    \rowcolor{GMue}München & 2\\\hline
    \rowcolor{GSch}Schwerin & 1\\\hline
    \end{tabular}}
    \end{minipage}};
\end{tikzpicture}\par\vspace{-.6mm}
\merkebox{Zeilen mit gleichem city-Wert bilden eine Gruppe. COUNT(*) zählt je Gruppe; pro Gruppe entsteht eine Ergebniszeile.}

\sectiontitle{2 \textperiodcentered{} Gruppen auswerten mit COUNT und AVG}
\begin{minipage}[t]{.43\linewidth}\vspace{0pt}
\workbox{Gib für männliche und weibliche Mitglieder jeweils das Geschlecht und die Anzahl aus.}
\begin{minipage}[t]{.62\linewidth}\vspace{0pt}
\smallcaption{SQL-Anweisung}
\sqlbox{\sqlkw{SELECT} gender, \sqlkw{COUNT}(*)\par\hspace{2mm}\sqlkw{AS} Anzahl\par\sqlkw{FROM} users\par\sqlkw{WHERE} gender = 'male'\par\hspace{2mm}\sqlkw{OR} gender = 'female'\par\sqlkw{GROUP BY} gender}
\end{minipage}\hfill
\begin{minipage}[t]{.35\linewidth}\vspace{0pt}
\smallcaption{Ergebnisrelation:\\2 Zeilen}
{\tabfont\begin{tabular}{|>{\ttfamily}l|>{\ttfamily\bfseries}r|}\hline
\rowcolor{Sky}gender & Anzahl\\\hline
female & 148\\\hline
male & 148\\\hline
\end{tabular}}
\end{minipage}
\end{minipage}\hfill
\begin{minipage}[t]{.55\linewidth}\vspace{0pt}
\workbox{Berechne für männliche und weibliche Mitglieder jeweils die durchschnittliche Körpergröße. Gib Geschlecht und Durchschnitt aus.}
\begin{minipage}[t]{.62\linewidth}\vspace{0pt}
\smallcaption{SQL-Anweisung}
\sqlbox{\sqlkw{SELECT} gender, \sqlkw{AVG}(centimeters)\par\hspace{2mm}\sqlkw{AS} Durchschnitt\par\sqlkw{FROM} users\par\sqlkw{WHERE} gender = 'male'\par\hspace{2mm}\sqlkw{OR} gender = 'female'\par\sqlkw{GROUP BY} gender}
\end{minipage}\hfill
\begin{minipage}[t]{.35\linewidth}\vspace{0pt}
\smallcaption{Ergebnisrelation: 2 Zeilen;\\Werte gekürzt}
{\tabfont\begin{tabular}{|>{\ttfamily}l|>{\ttfamily\bfseries}r|}\hline
\rowcolor{Sky}gender & Durchschnitt\\\hline
female & 167.52\ldots\\\hline
male & 178.39\ldots\\\hline
\end{tabular}}\par\vspace{.6mm}
\graynote{in cm; fehlende Größen\\zählen nicht mit}
\end{minipage}
\end{minipage}\par\vspace{.3mm}
\merkebox{COUNT(*) zählt die Zeilen einer Gruppe, AVG(spalte) bildet ihren Mittelwert.}

\sectiontitle{3 \textperiodcentered{} Gruppen filtern mit HAVING}
\begin{minipage}[t]{.43\linewidth}\vspace{0pt}
\smallcaption{Gruppen aus Abschnitt 1 mit HAVING COUNT(*) > 2}
\begin{tikzpicture}[x=1mm,y=1mm]
  \node[anchor=north west,inner sep=0] at (0,0) {%
    {\tabfont\begin{tabular}{|l|>{\ttfamily\bfseries}r|}\hline
    \rowcolor{Sky}city & COUNT(*)\\\hline
    \rowcolor{GBer}Berlin & 3\\\hline
    \rowcolor{GMue}München & 2\\\hline
    \rowcolor{GSch}Schwerin & 1\\\hline
    \end{tabular}}};
  \node[font=\fontsize{7.4}{8.6}\selectfont,text=Muted,anchor=north west] at (0,-17.4) {vor HAVING};
  \draw[bigarrow] (30,-9) -- (40,-9);
  \node[font=\fontsize{7.6}{9}\selectfont\ttfamily\bfseries,text=Blue,anchor=south] at (35,-7.6) {HAVING};
  \node[anchor=north west,inner sep=0] at (43,0) {%
    {\tabfont\begin{tabular}{|l|>{\ttfamily\bfseries}r|c}\cline{1-2}
    \rowcolor{Sky}city & COUNT(*) & \cellcolor{white}\\\cline{1-2}
    \rowcolor{GBer}Berlin & 3 & \cellcolor{white}\textcolor{Blue}{\ding{51}}\\\cline{1-2}
    \textcolor{Muted}{\sout{München}} & \textcolor{Muted}{2} & \textcolor{Muted}{\ding{55}}\\\cline{1-2}
    \textcolor{Muted}{\sout{Schwerin}} & \textcolor{Muted}{1} & \textcolor{Muted}{\ding{55}}\\\cline{1-2}
    \end{tabular}}};
  \node[font=\fontsize{7.4}{8.6}\selectfont,text=Muted,anchor=north west] at (43,-17.4) {nach HAVING};
\end{tikzpicture}\par\vspace{.8mm}
{\tabfont\begin{tabular}{|>{\ttfamily\bfseries}l|l|l|}\hline
\rowcolor{Sky}\normalfont Klausel & filtert \ldots & Zeitpunkt\\\hline
\textcolor{Blue}{WHERE} & einzelne Zeilen & vor GROUP BY\\\hline
\textcolor{Blue}{HAVING} & ganze Gruppen & nach GROUP BY\\\hline
\end{tabular}}
\end{minipage}\hfill
\begin{minipage}[t]{.55\linewidth}\vspace{0pt}
\workbox{Finde eingetragene Städte mit mehr als zwei Mitgliedern. Gib Stadt und Anzahl aus.}
\begin{minipage}[t]{.57\linewidth}\vspace{0pt}
\smallcaption{SQL-Anweisung}
\sqlbox{\sqlkw{SELECT} city, \sqlkw{COUNT}(*)\par\hspace{2mm}\sqlkw{AS} Anzahl\par\sqlkw{FROM} users\par\sqlkw{WHERE} city \sqlkw{IS NOT NULL}\par\sqlkw{GROUP BY} city\par\sqlkw{HAVING} \sqlkw{COUNT}(*) > 2}
\end{minipage}\hfill
\begin{minipage}[t]{.40\linewidth}\vspace{0pt}
\smallcaption{Ausschnitt: 3 von 26 Zeilen;\\Reihenfolge nicht festgelegt}
{\tabfont\begin{tabular}{|l|>{\ttfamily\bfseries}r|}\hline
\rowcolor{Sky}city & Anzahl\\\hline
\rowcolor{GBer}Berlin & 18\\\hline
Genf & 3\\\hline
\rowcolor{GMue}München & 24\\\hline
\ldots & \ldots\\\hline
\end{tabular}}\par\vspace{.6mm}
\graynote{Anzahlen aus der ganzen Tabelle;\\Genf: genau 3 > 2}
\end{minipage}
\end{minipage}\par\vspace{1.4mm}
{\fontsize{10}{12}\selectfont\textcolor{Blue}{\bfseries SELECT $\rightarrow$ FROM $\rightarrow$ WHERE $\rightarrow$ GROUP BY $\rightarrow$ HAVING $\rightarrow$ ORDER BY $\rightarrow$ LIMIT}}\par\vspace{.4mm}

\sectiontitle{4 \textperiodcentered{} Fotos pro Person zählen und sortieren}
\workbox{Gib die drei User-IDs mit den meisten Fotos und ihre Fotoanzahl aus. Sortiere nach Anzahl absteigend und bei Gleichstand nach User-ID aufsteigend.}
\begin{minipage}[t]{.42\linewidth}\vspace{0pt}
\smallcaption{SQL-Anweisung}
\sqlbox{\sqlkw{SELECT} user\_id, \sqlkw{COUNT}(*) \sqlkw{AS} Anzahl\par\sqlkw{FROM} photos\par\sqlkw{GROUP BY} user\_id\par\sqlkw{ORDER BY} Anzahl \sqlkw{DESC}, user\_id \sqlkw{ASC}\par\sqlkw{LIMIT} 3}
\end{minipage}\hfill
\begin{minipage}[t]{.19\linewidth}\vspace{0pt}
\smallcaption{Ergebnisrelation: 3 Zeilen}
{\tabfont\begin{tabular}{|>{\ttfamily}c|>{\ttfamily\bfseries}r|}\hline
\rowcolor{Sky}user\_id & Anzahl\\\hline
85 & 12\\\hline
187 & 12\\\hline
265 & 12\\\hline
\end{tabular}}\par\vspace{.6mm}
\graynote{Gleichstand $\rightarrow$\\user\_id aufsteigend}
\end{minipage}\hfill
\begin{minipage}[t]{.35\linewidth}\vspace{0pt}
\merkebox{GROUP BY garantiert keine Reihenfolge. Für eine Sortierung braucht es \mbox{ORDER BY}.}
\end{minipage}\par
\pagefooter{1}
\end{document}
